\section*{Hoofdstuk \ref{ch:g9:natural-risks} --- Leven met natuurrisico's}

\begin{solution}{exo:g9:natural-risks:1}
(a) Gevaar. (b) \omterm{def:g9:natural-risks:vulnerability}{Kwetsbaarheid}. (c) \omterm{def:g9:natural-risks:vulnerability}{Kwetsbaarheid}. (d) Gevaar. (e) Allebei:
de modderstromen zijn het gevaar; de stad op hun route bouwen maakt haar
kwetsbaar.
\end{solution}

\begin{solution}{exo:g9:natural-risks:2}
(a) Niet waar: waar niemand woont, kan niets schade oplopen, dus het \omterm{def:g9:natural-risks:risk}{risico}
is laag. (b) Waar. (c) Waar. (d) Niet waar: een waarschuwing verandert de
\omterm{def:g4:volcanoes-earthquakes-first-look:volcano}{vulkaan} niet; ze verlaagt de \omterm{def:g9:natural-risks:vulnerability}{kwetsbaarheid} van de mensen, die op tijd
kunnen vertrekken.
\end{solution}

\begin{solution}{exo:g9:natural-risks:3}
(a) Ja, dagen van tevoren. (b) Ja, uren tot dagen van tevoren. (c) Nee:
\omterm{def:g4:volcanoes-earthquakes-first-look:earthquake}{aardbevingen} kunnen niet worden voorspeld. (d) Ja: als de \omterm{def:g7:earthquakes:tsunami}{tsunami} eenmaal
is ontstaan, doet hij er uren over om de \omterm{def:g2:water-on-earth:ocean}{oceaan} over te steken, en de
waarschuwing gaat sneller.
\end{solution}

\begin{solution}{exo:g9:natural-risks:4}
(a) Het schudden van de grond: zwermen kleine \omterm{def:g4:volcanoes-earthquakes-first-look:earthquake}{aardbevingen} als \omterm{def:g7:volcanoes:magma}{magma}
opstijgt. (b) De helling van de flank: het zwellen van de \omterm{def:g4:volcanoes-earthquakes-first-look:volcano}{vulkaan}. (c) Hoe
de grond beweegt, met millimeters: het zwellen terwijl er \omterm{def:g7:volcanoes:magma}{magma} samenkomt.
(d) De gassen van de pluim: meer gas dat ontsnapt als \omterm{def:g7:volcanoes:magma}{magma} omhoogkomt.
\end{solution}

\begin{solution}{exo:g9:natural-risks:5}
(a) $45 \div 2 = \qty{22.5}{km/h}$. (b) $4 \times 2 = \qty{8}{km}$, ver
genoeg om hoger gelegen grond te bereiken. (c) Niemand zag de
\omterm{def:g4:volcanoes-earthquakes-first-look:eruption}{uitbarsting}, die door de \omterm{def:g5:natural-hazards:storm}{storm} werd verborgen, en niemand waarschuwde hen:
er was geen snelle manier om de informatie naar de stad te sturen, en geen
evacuatieplan. (d) Bijvoorbeeld: een evacuatieplan met routes naar hoger
gelegen grond, van tevoren geoefend; een waarschuwingsketen van de
wetenschappers naar de stad; meer instrumenten op de \omterm{def:g4:volcanoes-earthquakes-first-look:volcano}{vulkaan}.
\end{solution}

\begin{solution}{exo:g9:natural-risks:6}
(a) De \omterm{def:g5:water-shapes-the-land:flood}{overstroming} van \qty{10}{m} boven de \omterm{def:g1:landscapes-around-us:water}{rivier}. (b) Ongeveer
\qty{0.47}{km} voor de \omterm{def:g5:water-shapes-the-land:flood}{overstroming} van \qty{2}{m}, \qty{0.56}{km} voor
die van \qty{5}{m} en \qty{0.67}{km} voor die van \qty{10}{m} (aflezingen
tot op ongeveer \qty{0.05}{km} zijn goed). (c) Alleen de \omterm{def:g5:water-shapes-the-land:flood}{overstroming} van
\qty{10}{m}: de grond ligt hoger dan 2 en \qty{5}{m} boven de \omterm{def:g1:landscapes-around-us:water}{rivier}.
(d) Er is maar een kleine stijging van de \omterm{def:g1:landscapes-around-us:water}{rivier} nodig om haar te
bereiken, en kleine \omterm{def:g5:water-shapes-the-land:flood}{overstromingen} komen vaker voor dan grote.
\end{solution}

\begin{solution}{exo:g9:natural-risks:7}
(a) Groot gevaar, kleine \omterm{def:g9:natural-risks:vulnerability}{kwetsbaarheid}: matig \omterm{def:g9:natural-risks:risk}{risico}. (b) Klein gevaar,
grote \omterm{def:g9:natural-risks:vulnerability}{kwetsbaarheid}: matig \omterm{def:g9:natural-risks:risk}{risico}. (c) Groot gevaar, grote \omterm{def:g9:natural-risks:vulnerability}{kwetsbaarheid}:
hoog \omterm{def:g9:natural-risks:risk}{risico}. (d) Klein gevaar, kleine \omterm{def:g9:natural-risks:vulnerability}{kwetsbaarheid}: laag \omterm{def:g9:natural-risks:risk}{risico}.
\end{solution}

\begin{solution}{exo:g9:natural-risks:8}
(a) C, dan B, dan A. (b) De schoren vormen driehoeken, die niet kunnen
overhellen zonder een schoor uit te rekken of in te drukken: het skelet
blijft in vorm in plaats van over te hellen. (c) De isolatoren geven mee
als de grond eronder heen en weer beweegt, zodat veel minder van de
beweging de school bereikt.
\end{solution}

\begin{solution}{exo:g9:natural-risks:9}
(a) $1 \div 100 = 1~\%$ per jaar. (b) De kans is elk jaar dezelfde, wat er
het jaar ervoor ook gebeurde: de \omterm{def:g5:water-shapes-the-land:flood}{overstroming} kan volgend jaar weer komen.
(c) $300 \div 100 = 3$ \omterm{def:g5:water-shapes-the-land:flood}{overstromingen}, gemiddeld.
\end{solution}

\begin{solution}{exo:g9:natural-risks:10}
(a) $3600 \div 15 = 240$ metingen. (b) $240 \times 24 \times 365 =
\num{2102400} \approx 2.1 \times 10^6$ metingen. (c) $2400 \div 800 =
\qty{3}{h}$. (d) In die drie uur kunnen de waarschuwingscentra de \omterm{def:g1:landscapes-around-us:water}{kust}
waarschuwen en kunnen de mensen naar hoger gelegen grond gaan.
\end{solution}

\begin{solution}{exo:g9:natural-risks:11}
(a) $4 \times 3 \times 3.8 = \qty{45.6}{L}$. (b) $45.6 \div 1.5 = 30.4$,
dus 31 flessen. (c) Bijvoorbeeld: houdbaar voedsel, om te eten als de
winkels dicht zijn; een zaklamp, als de stroom uitvalt; een radio op
batterijen of met een slinger, om de waarschuwingen te horen; een
EHBO-doos, om kleine verwondingen te verzorgen; warme dekens. (d) Zodat
het meteen kan worden meegenomen als het gezin snel weg moet.
\end{solution}

\begin{solution}{exo:g9:natural-risks:12}
Bijvoorbeeld: (1) een \omterm{def:g9:natural-risks:hazard-map}{gevarenkaart} van de modderstromen tekenen, \omterm{def:g1:landscapes-around-us:landform}{dal} per
\omterm{def:g1:landscapes-around-us:landform}{dal} (kant van het gevaar); (2) \omterm{def:g7:earthquakes:seismometer}{seismometers}, hellingmeters,
gps-ontvangers en gasmeters op de \omterm{def:g4:volcanoes-earthquakes-first-look:volcano}{vulkaan} plaatsen (kant van het gevaar);
(3) een snelle waarschuwingsketen van het observatorium naar de stad
opzetten, met sirenes (kant van de \omterm{def:g9:natural-risks:vulnerability}{kwetsbaarheid}); (4) routes naar hoger
gelegen grond plannen en aangeven, en ze met oefeningen oefenen (kant van
de \omterm{def:g9:natural-risks:vulnerability}{kwetsbaarheid}); (5) geen nieuwe woningen of scholen bouwen op de route
van de modderstromen (kant van de \omterm{def:g9:natural-risks:vulnerability}{kwetsbaarheid}).
\end{solution}

\begin{solution}{pb:g9:natural-risks:1}
\textbf{1.} De \omterm{def:g1:landscapes-around-us:water}{rivier} staat op \qty{143}{m} bij C, \qty{137}{m} bij A en
\qty{126}{m} bij B; op de \omterm{def:g2:drawing-to-map:map}{kaart} ligt C in het westen en B in het oosten,
dus de \omterm{def:g1:landscapes-around-us:water}{rivier} stroomt naar het oosten.

\textbf{2.} A: $137 - 137 = \qty{0}{m}$; B: $276 - 126 = \qty{150}{m}$;
C: $236 - 143 = \qty{93}{m}$.

\textbf{3.} Alleen A, dat minder dan \qty{5}{m} boven de \omterm{def:g1:landscapes-around-us:water}{rivier} ligt.

\textbf{4.} Ongeveer \qty{0.57}{km} (tussen 0.5 en \qty{0.65}{km}).

\textbf{5.} De stijging over \qty{100}{m} is de helft van de stijging over
\qty{200}{m}: A $1.4 \div 2 = 0.7~\%$; B $83 \div 2 = 41.5~\%$;
C $20 \div 2 = 10~\%$.

\textbf{6.} B: de helling is er heel steil, en los \omterm{def:g3:stones-pebbles-sand:rock}{gesteente} en losse
grond kunnen langs een steile helling naar beneden glijden, vooral na
zware \omterm{def:g4:clouds-rain-snow:precipitation}{regen}.

\textbf{7.} A: ongeveer \qty{1}{km}; B: bijna op de \omterm{def:g7:earthquakes:fault}{breuk} (minder dan
\qty{0.05}{km}); C: ongeveer \qty{1.3}{km}.

\textbf{8.} Als de \omterm{def:g7:earthquakes:fault}{breuk} verschuift, breekt en verschuift de grond langs
de \omterm{def:g7:earthquakes:fault}{breuk}, wat een gebouw uit elkaar kan scheuren; en een \omterm{def:g4:volcanoes-earthquakes-first-look:earthquake}{aardbeving} op de
\omterm{def:g7:earthquakes:fault}{breuk} zou de grond daar heel hard laten schudden.

\textbf{9.}

\begin{center}
{\small
\begin{tabular}{|c|c|c|c|}
\hline
plaats & in de overstromingszone? & steile helling? & op de \omterm{def:g7:earthquakes:fault}{breuk}? \\
\hline
A & ja & nee & nee \\
B & nee & ja & ja \\
C & nee & nee & nee \\
\hline
\end{tabular}}
\end{center}

\textbf{10.} Plaats C: het is de enige plaats met drie keer ``nee'':
\qty{93}{m} boven de \omterm{def:g1:landscapes-around-us:water}{rivier}, een flauwe helling van \qty{10}{\%}, meer
dan \qty{1}{km} van de \omterm{def:g7:earthquakes:fault}{breuk}.

\textbf{11.} Bijvoorbeeld: diagonale schoren en schijfwanden in het
skelet, stevige verbindingen, of isolatoren onder het gebouw; de
leerlingen oefenen voor \omterm{def:g4:volcanoes-earthquakes-first-look:earthquake}{aardbevingen} (bukken, dekken, vasthouden).

\textbf{12.} \textbf{17 huizen} van de 28 staan in de overstromingszone:
$17 \div 28 \approx 0.61$, ongeveer 3 huizen op 5.
\end{solution}
